\newpage
\section{Conclusiones}

En este trabajo se desarrolla un sistema híbrido de control y automatización para una grúa
portacontenedores tipo STS, organizado en una arquitectura jerárquica de tres niveles que
coordina la traslación del carro y el izaje de la carga de forma semi-automática, simultánea y
segura. La estrategia combina un control de eventos discretos con un control regulatorio en
tiempo discreto, integrando las capas de desempeño y de seguridad en un único modelo.

El control regulatorio (Nivel 2) cumple el seguimiento de las consignas de velocidad de ambos
ejes mediante lazos PI sintonizados por asignación de polos, complementados con observadores de
estado para estimar la velocidad a partir de la sola medición de posición y con una compensación
de gravedad que descarga el término integral del izaje. El control de balanceo (\textit{anti-sway})
desacoplado amortigua eficazmente la oscilación pendular de la carga, transformando la cinemática
del carro en un amortiguamiento viscoso sin depender de la aproximación de ángulos pequeños.

El control supervisor (Nivel 1) se implementa como un autómata de procesos paralelos que coordina
los modos de operación (reposo, manual y automático), gestiona la lógica de toma y liberación de
la carga mediante los \textit{twistlocks}, releva el perfil de obstáculos con el LIDAR y planifica
trayectorias con evasión de colisiones. Sus límites suaves de velocidad y el suavizado de
consignas garantizan movimientos graduales y dentro del área de trabajo válida, evitando que la
operación alcance los límites físicos de forma abrupta.

El control de protección (Nivel 0) constituye la capa de mayor autoridad: supervisa de forma
permanente la planta y, ante cualquier anomalía, sobrescribe a los niveles superiores aplicando
los frenos mecánicos. Su frenado selectivo detiene únicamente el eje afectado, mientras que el
monitor de comunicaciones (\textit{watchdog}) detecta la pérdida de enlace con el supervisor y
fuerza la parada de emergencia, garantizando que la grúa nunca quede operando con un comando
desactualizado.

El comportamiento del sistema se valida satisfactoriamente mediante simulación, apoyada en una
animación 2D en tiempo real que permite observar la operación en los distintos modos. Asimismo, la
lógica de control se traslada a un PLC virtual programado en CODESYS (combinando SFC y texto
estructurado) y se comunica con la simulación de la planta a través del protocolo OPC-UA,
comprobándose que el controlador implementado responde de igual manera que el modelo original.

Como trabajos futuros se propone llevar el control a hardware real sobre un PLC físico, sustituir
la comunicación OPC-UA por un enlace de menor latencia apto para control en tiempo real, y
profundizar en las estrategias de planificación de trayectorias y de control anti-balanceo para
optimizar los tiempos de ciclo de la grúa.
